\exercisesheader{}

% 27

\eoce{\qt{Melempar dadu\label{roll_die}} Hitung peluang-peluang
berikut dan tentukan model distribusi peluang yang tepat untuk setiap
kasus. Anda melempar sebuah dadu seimbang 5 kali. Berapa peluang memperoleh
\begin{parts}
\item angka 6 pertama pada lemparan kelima?
\item tepat tiga angka 6?
\item angka 6 ketiga pada lemparan kelima?
\end{parts}
}{}

% 28

\eoce{\qt{Bermain dart\label{play_darts}} Hitung peluang-peluang berikut
dan tentukan model distribusi peluang yang tepat untuk setiap kasus.
Seorang pemain dart yang sangat terampil dapat mengenai pusat sasaran
(lingkaran merah di tengah papan dart) dengan peluang 65\%. Berapa
peluang ia
\begin{parts}
\item mengenai pusat sasaran untuk kali ke-$10$ pada percobaan ke-$15$?
\item mengenai pusat sasaran 10 kali dalam 15 percobaan?
\item mengenai pusat sasaran pertama kali pada percobaan ketiga?
\end{parts}
}{}

% 29

\eoce{\qt{Pengambilan sampel di sekolah\label{sampling_at_school}} Untuk proyek
mata kuliah sosiologi, Anda diminta melakukan survei terhadap 20 siswa
di sekolah Anda. Anda
memutuskan untuk berdiri di luar kafetaria asrama pada suatu malam dan melakukan
survei terhadap sampel acak 20 siswa yang keluar dari kafetaria seusai makan malam.
Penghuni asrama Anda terdiri atas 45\% laki-laki dan 55\% perempuan.
\begin{parts}
\item Model peluang mana yang paling tepat untuk menghitung
peluang bahwa responden ke-$4$ merupakan responden perempuan ke-$2$?
Jelaskan.
\item Hitung peluang pada bagian (a).
\item Berikut adalah tiga skenario ketika responden ke-$4$
merupakan responden perempuan ke-$2$ (dengan $M$ menyatakan laki-laki dan $F$ menyatakan perempuan):
\[ \{M, M, F, F\}, \{M, F, M, F\}, \{F, M, M, F\} \]
Kesamaan di antara skenario-skenario ini adalah bahwa responden pada percobaan terakhir selalu
perempuan. Dalam tiga percobaan pertama terdapat 2 laki-laki dan 1 perempuan. Gunakan
koefisien binomial untuk mengonfirmasi bahwa ada 3 cara mengurutkan 2 laki-laki dan
1 perempuan.
\item Gunakan temuan pada bagian (c) untuk menjelaskan mengapa koefisien
dalam rumus distribusi binomial negatif adalah ${n-1 \choose k-1}$, sedangkan
koefisien binomial adalah ${n \choose k}$.
\end{parts}
}{}

% 30

\eoce{\qt{Servis dalam bola voli\label{serving_volleyball}} Seorang pemain
bola voli yang kurang terampil memiliki peluang 15\% untuk melakukan servis yang masuk, yaitu
memukul bola agar melintasi net dengan lintasan sedemikian rupa sehingga bola akan
jatuh di lapangan tim lawan. Misalkan servis-servisnya saling independen.
\begin{parts}
\item Berapa peluang bahwa pada percobaan ke-$10$ ia akan melakukan
servis suksesnya yang ke-$3$?
\item Misalkan ia telah melakukan dua servis sukses dalam sembilan percobaan. Berapa
peluang bahwa servisnya yang ke-$10$ akan sukses?
\item Meskipun bagian (a) dan (b) membahas skenario yang sama,
peluang yang Anda hitung seharusnya berbeda. Dapatkah Anda menjelaskan penyebab
perbedaan ini?
\end{parts}
}{}
